\section{A polynomial sample lower bound at each fixed critical depth}
\label{sec:critical-polynomial-sample-lower}

We can strengthen the consistent-dictator argument by constructing a
low-rate interpolant on a small coordinate projection. This remains a
polynomial lower bound, rather than an exponential lower bound in parity
degree, and the width limit is taken before the dimension limit.

\begin{lemma}[Amplifying a first-feature mean through its second moment]
\label{lem:critical-mixture-mean-seed}
Let $\mu_1$ be the uncontrolled first-feature law. Suppose a feature law
$\nu$ satisfies $D(\nu\Vert\mu_1)\le D$, where $D\ge1$, and its mean
$h$ satisfies $g(x)h(x)\ge a$ at every input, with $0<a\le1$ and
$g$ Boolean. Set $R=D/a^2$ and $q=R^{1/L}$. If $q\le D$, then
\[
 J_L(g)\le5000L R^{1/L}.
\]
\end{lemma}

\begin{proof}
Put $\theta=q/D\le1$. The first controlled law
$(1-\theta)\mu_1+\theta\nu$ has entropy at most $\theta D=q$ by
convexity, and its second moment dominates
$\theta hh^T$. Thus it has a tracked rank-one vector
$b_1=\sqrt\theta h$ with
\[
 g(x)b_1(x)\ge\sqrt\theta a=q^{-(L-1)/2}.
\]
This is a second-moment component; it need not equal the first
controlled law's mean.

For $L=2$, choose the readout $z=h/(2a)$. Its cost is at most
$1/(8\theta a^2)=q/8$, proving the assertion in this case.
Suppose $L\ge3$ and put $c=1/100$. For every $0\le v\le1$ and
$u\ge0$, the elementary estimate used in the constant-target
construction gives
\[
 \mathbb E\tanh(u+\sqrt v Z)\ge c\min\{u,1\}.
\]
At the second hidden layer shift the reference Gaussian by
$\sqrt q\,b_1$. Its mean-shift entropy is at most $q/2$, because
$Q_1\succeq b_1b_1^T$. After tanh, its actual mean $b_2$ satisfies
$g(x)b_2(x)\ge c q^{-(L-2)/2}$. Its second moment dominates
$b_2b_2^T$.

At each hidden layer $s=3,\ldots,L-1$, shift the reference Gaussian
by $(\sqrt q/c)b_{s-1}$. The entropy is at most $q/(2c^2)$, and
the new feature mean satisfies
\[
 g(x)b_s(x)\ge c q^{-(L-s)/2}.
\]
The induction uses only that every reference diagonal is at most one,
not that the diagonals or feature magnitudes are equal. If some shift
coordinate exceeds one in magnitude, the displayed minimum estimate
still gives the required lower bound. All shift vectors lie in the
corresponding covariance range; the formulas therefore also apply to
singular covariances. Tanh pushforward cannot increase the entropy.

Finally take $z=b_{L-1}/(2c q^{-1/2})$. It has the required global
margin and Gaussian readout cost at most $q/(8c^2)$. Summing all
costs gives the claimed bound.
\end{proof}

\begin{lemma}[A low-cost first-feature law for a coordinate memorizer]
\label{lem:critical-projected-memorizer}
Suppose $n$ labeled cube inputs have distinct projections onto $r\ge1$
specified coordinates. There is a Boolean function $g$ on those
coordinates agreeing with all labels and a first-feature law $\nu$
whose mean $h$ satisfies $gh\ge a$, where
\[
 b=16\log(32(n+1)),\qquad
 D=d b^2(r+1)^2,\qquad a=\frac1{4(n+1)},
 \qquad D(\nu\Vert\mu_1)\le D.
\]
\end{lemma}

\begin{proof}
Let $p\le n$ be the number of positive labeled projections. Define $g$
to be positive exactly at those projections and negative elsewhere.
For a positive projected point $y^j\in\{-1,1\}^r$, use the Gaussian
first-row law with its mean shifted to the affine form
\[
 A_j(x)=b\left(\sum_{i=1}^r y_i^j x_i-r+1\right).
\]
Its deterministic part is at least $b$ at $y^j$ and at most $-b$
elsewhere. At every input its remaining Gaussian variance is
$(d+1)/d\le2$. If $s_j(x)$ is its mean tanh feature, Gaussian tails
and $1-\tanh t\le2e^{-2t}$ give
\[
 |s_j(x)-\operatorname{sign}A_j(x)|
 \le2e^{-b}+4e^{-b^2/16}
 \le\frac1{16(n+1)}.
\]
The same bound holds for the constant-positive feature mean $\tau$
obtained by shifting only the bias to $b$.

For $p\ge1$, put
\[
 F(x)=(p-1)\tau+\sum_{j=1}^p s_j(x),\qquad A=2p-1.
\]
Replacing each mean feature by its ideal sign gives exactly $g(x)$:
at a positive projected point exactly one point indicator is positive.
The total approximation error is at most
$A/[16(n+1)]<1/8$. Choose $\nu$ as the mixture of the corresponding
shifted first-row Gaussian laws, assigning each point indicator weight
$1/A$ and the constant law weight $(p-1)/A$, and push it through tanh.
Then its mean is $h=F/A$, so $gh\ge7/(8A)\ge a$.
For $p=0$, take only the constant-negative shifted law; it satisfies
the same conclusion.

Every shift vector has squared Euclidean norm at most
$b^2(r+1)^2$. The reference first-row law is
$N(0,I_{d+1}/d)$, so each shifted Gaussian costs at most
$d b^2(r+1)^2/2$. Convexity of entropy and data processing prove
the asserted bound, including all unused input coordinates.
\end{proof}

\begin{theorem}[Critical parity sample lower bound]
\label{thm:critical-polynomial-sample-lower}
Fix $L\ge2$ and $0<\eta<1$. Let the target be parity on $k\le d$
specified input coordinates, and put
\[
 M=\min\{k,d^{1/(L-1)}\}.
\]
Suppose $M\to\infty$ and $n\le M^{(1-\eta)/2}$. Draw $n$ independent
uniform cube inputs labeled by this parity. For all sufficiently large
$M$, with training-sample probability at least $1-1/(2M)$, the globally
margin-conditioned posterior at
$\gamma=1$ has Gibbs error tending to $1/2$ as $N\to\infty$ at
this fixed $d,L,n$. If $p_N(x)$ is its posterior probability of the
correct label, then on the same event
\[
 \liminf_{N\to\infty}\Pr_x\{p_N(x)<2/3\}\ge1/4.
\]
For full parity, this gives $n\le d^{(1-\eta)/(2(L-1))}$. If
$d\ge k^{L-1}$, it gives $n\le k^{(1-\eta)/2}$.
\end{theorem}

\begin{proof}
Use the first
\[
 r=\left\lceil\log_2\bigl(M(n+1)^2\bigr)\right\rceil
\]
input coordinates. At the specified sample sizes $r=O(\log M)<d$
for all sufficiently large $M$, because $M\le d$. A union bound over pairs gives
projection-collision probability at most
$n(n-1)2^{-r}/2\le1/(2M)$.

On the complementary event, the memorizer lemma and the amplification
lemma give a training-consistent function of rate at most
\[
 5000L\left[16d b^2(r+1)^2(n+1)^2\right]^{1/L}.
\]
The amplification hypothesis holds for large $M$: it is equivalent
to $D^{L-1}\ge16(n+1)^2$, where $D=d b^2(r+1)^2$.
At our sample sizes this follows from $M\le d^{1/(L-1)}$ and
$n^2\le M^{1-\eta}$ (including $L=2$, where the positive power
slack is at least $M^\eta$).

Any Boolean function depending on every one of the $k$ parity coordinates
has at least $k$ relevant augmented coordinates. By
Theorem~\ref{thm:critical-fixed-depth-ordering}, its critical rate is at least
\[
 \frac1{64}\min\{d^{1/(L-1)},(dk)^{1/L}\}
       =\frac1{64}(dM)^{1/L}.
\]
The ratio of the
preceding competitor upper bound to this lower bound is at most
\[
 C L\,M^{-\eta/L}[\log(M+1)]^{4/L},
\]
which tends to zero. Thus there is a positive rate gap between a
training-consistent competitor and every function retaining all
parity coordinates.

The exact finite-domain posterior-rate principle now implies that
posterior mass on the latter functions tends to zero in the width
limit. A function omitting any one parity coordinate has uniform
parity error exactly $1/2$, even if it depends on irrelevant coordinates.
This proves the Gibbs limit. Finally
$\mathbb E_xp_N(x)\to1/2$ and Markov's inequality give
$\Pr_x\{p_N(x)<2/3\}\ge1-(3/2)\mathbb E_xp_N(x)$.
\end{proof}

The theorem concerns the actual posterior, but takes width to infinity
at fixed dimension, depth, and sample size. It gives no finite-width
threshold uniform in growing dimension and no corresponding lower
bound for posterior-majority prediction at confidence threshold $1/2$.

\begin{corollary}[Slowly growing learner depth]
\label{cor:critical-growing-depth-sample-lower}
Fix $0<\eta<1$. For full parity, the same conclusion holds at sample sizes
$n\le d^{(1-\eta)/(2(L-1))}$ along integer depths $L=L(d)\ge2$ satisfying
\[
 L=o\!\left(\sqrt{\frac{\log d}{\log\log d}}\right).
\]
For each $d$, the width limit is taken with its corresponding $L(d)$
and sample size fixed.
\end{corollary}

\begin{proof}
Here $M=d^{1/(L-1)}\to\infty$. The projection and Gaussian-mixture bounds
are uniform in $L\ge2$:
$r\le2\log_2d+O(1)$ and $b\le C\log(d+1)$ at the specified sample
sizes. The preceding rate ratio is therefore at most
\[
 C L\,d^{-\eta/[L(L-1)]}[\log(d+1)]^{4/L}
\]
with universal $C$. Its logarithm tends to minus infinity under the
stated depth assumption. The mixture hypothesis and collision bound
also hold uniformly, so the same rate-gap argument applies.
\end{proof}
